% Draft abstract based on the toy-example results; revise once the industrial
% comparison is done.
Advanced process controllers based on learned models, such as TD-MPC, can
perform well but are hard to inspect, verify and audit. This thesis
investigates an alternative in which a large language model (LLM) is used
\emph{offline} as a meta-supervisor: it reads logs, scores and decision
traces from a simulated process and writes a small, deterministic supervisory
function in Python. The function runs on top of fixed PID loops, raises
per-loop anomaly flags and adjusts setpoints. Every candidate passes a static
security check and a sandboxed evaluation on a battery of fault scenarios, and
is only promoted if it improves on the current champion without regressing
any single scenario.

The approach is developed on three toy processes of increasing difficulty: a
single leaking tank, a two-tank cascade, and the quadruple-tank process. On
all three, the generated supervisors reduced the score on the development
scenarios by 72--96\% compared with PID control alone. On the quadruple-tank process, supervisors
generated with LLM reasoning enabled matched a perfect-detection oracle on
both the development and the held-out scenarios within two to three
generations, whereas the same search without reasoning stagnated at about
4.5 times the oracle score. The experiments also show that the design of the
evaluation, including safeguards against specification gaming and checks that
the test bed itself is sound, matters as much as the choice of model.
\TODO{Add the result of the comparison on the industrial process and against TD-MPC.}
